% !TEX program = pdflatex
%=====================================================================
%  Criptografía y Seguridad - ITBA
%  Clase 10 — Seguridad en la empresa
%  Apunte integral: teoría + práctica
%=====================================================================
\documentclass[11pt,a4paper]{article}

%---------------------------------------------------------------------
%  Paquetes
%---------------------------------------------------------------------
\usepackage[utf8]{inputenc}
\usepackage[T1]{fontenc}
\usepackage[spanish,es-tabla,es-noshorthands]{babel}
\usepackage{lmodern}
\usepackage{microtype}

\usepackage{amsmath,amssymb,amsthm}
\usepackage{mathtools}
\usepackage{geometry}
\geometry{a4paper,margin=2.4cm,headheight=22pt}

\usepackage{xcolor}
\usepackage{graphicx}
\usepackage{enumitem}
\usepackage{booktabs}
\usepackage{array}
\usepackage{tcolorbox}
\tcbuselibrary{skins,breakable}
\usepackage{fancyhdr}
\usepackage{titlesec}
\usepackage{hyperref}
\usepackage{listings}
\usepackage{caption}

\usepackage{tikz}
\usetikzlibrary{shapes.geometric,shapes.misc,shapes.symbols,arrows.meta,positioning,calc,
                fit,backgrounds,decorations.pathmorphing,shadows.blur,chains,mindmap,trees}

%---------------------------------------------------------------------
%  Paleta de color (inspirada en las diapositivas de la cátedra)
%---------------------------------------------------------------------
\definecolor{itbaCyan}{HTML}{6EC6D9}
\definecolor{itbaCyanDark}{HTML}{2E8B9E}
\definecolor{itbaBlue}{HTML}{AEDCEA}
\definecolor{boxBlue}{HTML}{BBD4F0}
\definecolor{slideGray}{HTML}{ECECEC}
\definecolor{practGreen}{HTML}{4F8A3D}
\definecolor{practGreenL}{HTML}{E2EFD9}
\definecolor{attackRed}{HTML}{C0392B}
\definecolor{codeBg}{HTML}{F5F5F5}
\definecolor{purpleZTA}{HTML}{7B5BA6}
\definecolor{purpleZTAL}{HTML}{ECE3F5}
\definecolor{riskOrange}{HTML}{E67E22}

%---------------------------------------------------------------------
%  Hipervínculos
%---------------------------------------------------------------------
\hypersetup{
  colorlinks=true,
  linkcolor=itbaCyanDark,
  urlcolor=itbaCyanDark,
  citecolor=itbaCyanDark,
  pdftitle={Clase 10 - Seguridad en la empresa},
  pdfauthor={Criptografia y Seguridad - ITBA}
}

%---------------------------------------------------------------------
%  Encabezados y pies
%---------------------------------------------------------------------
\pagestyle{fancy}
\fancyhf{}
\lhead{\small\textsc{Criptografía y Seguridad — ITBA}}
\rhead{\small Clase 10 · Seguridad en la empresa}
\cfoot{\small\thepage}
\renewcommand{\headrulewidth}{0.4pt}
\renewcommand{\headrule}{\hbox to\headwidth{\color{itbaCyanDark}\leaders\hrule height \headrulewidth\hfill}}

%---------------------------------------------------------------------
%  Títulos de sección con barra de color (estilo diapositiva)
%---------------------------------------------------------------------
\titleformat{\section}
  {\normalfont\Large\bfseries\color{itbaCyanDark}}
  {\thesection}{0.6em}{}
  [\vspace{-0.4em}{\color{itbaCyan}\titlerule[2pt]}]
\titleformat{\subsection}
  {\normalfont\large\bfseries\color{itbaCyanDark}}
  {\thesubsection}{0.6em}{}
\titleformat{\subsubsection}
  {\normalfont\normalsize\bfseries\color{black}}
  {\thesubsubsection}{0.6em}{}

%---------------------------------------------------------------------
%  Cajas reutilizables
%---------------------------------------------------------------------
% Caja "diapositiva" (recuadro gris de las slides)
\newtcolorbox{slidebox}[1][]{
  enhanced, colback=slideGray, colframe=black!55, boxrule=0.5pt,
  arc=1pt, left=6pt,right=6pt,top=4pt,bottom=4pt, #1}

% Caja de definición
\newtcolorbox{defbox}[1]{
  enhanced, breakable, colback=itbaBlue!25, colframe=itbaCyanDark,
  boxrule=1pt, arc=2pt, left=8pt,right=8pt,top=6pt,bottom=6pt,
  fonttitle=\bfseries, coltitle=white, title={#1}}

% Caja de nota práctica
\newtcolorbox{practbox}[1]{
  enhanced, breakable, colback=practGreenL, colframe=practGreen,
  boxrule=1pt, arc=2pt, left=8pt,right=8pt,top=6pt,bottom=6pt,
  fonttitle=\bfseries, coltitle=white, title={#1}}

% Caja de atención / ataque
\newtcolorbox{warnbox}[1]{
  enhanced, breakable, colback=attackRed!8, colframe=attackRed,
  boxrule=1pt, arc=2pt, left=8pt,right=8pt,top=6pt,bottom=6pt,
  fonttitle=\bfseries, coltitle=white, title={#1}}

% Caja anécdota / dicho en clase
\newtcolorbox{anecbox}[1]{
  enhanced, breakable, colback=riskOrange!8, colframe=riskOrange,
  boxrule=1pt, arc=2pt, left=8pt,right=8pt,top=6pt,bottom=6pt,
  fonttitle=\bfseries, coltitle=white, title={#1}}

%---------------------------------------------------------------------
%  Estilo de código
%---------------------------------------------------------------------
\lstdefinestyle{hcl}{
  basicstyle=\ttfamily\footnotesize,
  backgroundcolor=\color{codeBg},
  keywordstyle=\color{itbaCyanDark}\bfseries,
  commentstyle=\color{practGreen}\itshape,
  stringstyle=\color{attackRed},
  breaklines=true, frame=single, rulecolor=\color{black!30},
  showstringspaces=false, columns=fullflexible, tabsize=2,
  xleftmargin=10pt, framexleftmargin=8pt}

\newcommand{\factor}[1]{\textcolor{itbaCyanDark}{\textbf{#1}}}

% Cita textual del profesor
\newcommand{\prof}[1]{\par\medskip\noindent\textit{``#1''}\\[-2pt]\hspace*{1em}%
  {\footnotesize\textcolor{black!60}{--- dicho en clase}}\par\medskip\noindent}

%=====================================================================
\begin{document}
%=====================================================================

%---------------------------------------------------------------------
%  Portada
%---------------------------------------------------------------------
\begin{titlepage}
\centering
\vspace*{1.0cm}

\begin{tikzpicture}
  \node[rounded corners=2pt, fill=itbaCyan, text=black, inner sep=14pt,
        minimum width=0.8\textwidth, align=center] (t)
        {\Huge\bfseries Criptografía y Seguridad};
\end{tikzpicture}

\vspace{0.6cm}

{\LARGE Seguridad en aplicaciones}\\[0.25cm]
{\Huge\bfseries\color{itbaCyanDark} Seguridad en la empresa}

\vspace{0.8cm}

% Ilustración: bóveda / caja fuerte (recreación del icono de portada)
\begin{tikzpicture}[scale=1.0]
  \fill[black!12] (-3.2,-2.4) rectangle (3.2,2.8);
  \draw[line width=2pt, black!55] (-3.2,-2.4) rectangle (3.2,2.8);
  \fill[itbaCyanDark!75] (-2.4,-2.0) rectangle (2.4,2.4);
  \draw[line width=3pt, black!70] (-2.4,-2.0) rectangle (2.4,2.4);
  \draw[line width=1.2pt, itbaCyan] (-2.0,-1.6) rectangle (2.0,2.0);
  \draw[line width=2pt, black!70] (0,-2.0) -- (0,2.4);
  \foreach \a in {0,45,...,315}{
    \draw[line width=2pt, black!70] (-1.1,0.2) -- ++(\a:0.6);
  }
  \draw[line width=2.5pt, black!70] (-1.1,0.2) circle (0.42);
  \fill[itbaCyan] (-1.1,0.2) circle (0.16);
  \draw[line width=2.5pt, black!70] (1.1,0.2) circle (0.42);
  \fill[itbaCyan] (1.1,0.2) circle (0.16);
  \foreach \y in {-1.2,-0.4,0.8,1.6}{
     \fill[black!70] (-2.2,\y) circle (0.07);
     \fill[black!70] (2.2,\y) circle (0.07);
  }
\end{tikzpicture}

\vfill
{\large Apunte integral — Teoría + Práctica}\\[0.2cm]
{\normalsize Clase 10 \quad·\quad Lectura recomendada: \emph{NIST SP 800-207} (caps.\ 2 y 3)}\\[0.4cm]
{\small Instituto Tecnológico de Buenos Aires (ITBA)}

\vspace{0.6cm}
{\footnotesize\itshape Documento elaborado a partir de las transcripciones de la clase
teórica (Prof.\ Leandro Suárez) y de las diapositivas oficiales de la teórica y la práctica.
Esta clase \textbf{no está en Bishop}: está enfocada en el panorama profesional contemporáneo de
la seguridad en organizaciones.}
\end{titlepage}

\tableofcontents
\newpage

%=====================================================================
\section{Introducción y contexto}
%=====================================================================

Esta clase cierra el bloque de \textbf{Seguridad en aplicaciones} pasando de la teoría
(autenticación, control de acceso, principios de diseño, vulnerabilidades) a la
\textbf{visión profesional}: cómo una empresa real adopta distintas estrategias para
\emph{protegerse} de un abanico amplio de amenazas, equilibrando seguridad y negocio.

\begin{anecbox}{Dónde se sitúa esta clase}
\prof{Esto de la clase que vamos a ver hoy no está en Bishop. [\dots] está más enfocado en
lo que es el panorama contemporáneo, lo profesional: cómo se maneja esto en un ambiente
profesional de una organización, una corporación, incluso una Pyme con buenas prácticas.}
La lectura recomendada para profundizar es el estándar \textbf{NIST SP 800-207} (Zero Trust),
capítulos 2 y 3.
\end{anecbox}

\begin{defbox}{Tríada CIA — el marco de toda la seguridad (recordatorio)}
Toda política de seguridad busca \emph{preservar} tres propiedades sobre la información:
\begin{itemize}[leftmargin=1.4em,itemsep=2pt]
  \item \textbf{Confidencialidad}: guardar la info en secreto; acceso sólo a autorizados.
  \item \textbf{Integridad}: confiabilidad de los datos. Se distingue \emph{integridad de
        datos} (no se alteran) e \emph{integridad de origen} (provienen de quien dicen).
  \item \textbf{Disponibilidad}: poder usar la información cuando se la necesita.
\end{itemize}
El propósito de todo lo que veremos hoy es, en el fondo, \textbf{forzar el cumplimiento de la
tríada CIA dentro de la organización}, tanto de forma \textbf{proactiva} (barreras, controles,
evitar exponer vulnerabilidades) como \textbf{reactiva} (qué medidas tomar una vez que ya nos
atacaron).
\end{defbox}

\subsection{Una mirada del temario desde la práctica: la certificación CISSP}

La clase práctica abre presentando la certificación profesional de referencia en este campo:
\textbf{CISSP} (\emph{Certified Information System Security Professional}, de ISC2). Sus
\textbf{8 dominios de conocimiento} son un excelente mapa mental de \emph{toda} la materia de
seguridad de la empresa:

\begin{practbox}{Los 8 dominios de CISSP (aporte de la práctica)}
\begin{enumerate}[leftmargin=1.6em,itemsep=1pt]
  \item Gestión de Seguridad y riesgos
  \item Seguridad de activos
  \item Arquitectura e ingeniería de la seguridad
  \item Seguridad de comunicaciones y redes
  \item Gestión de identidades y accesos (IAM)
  \item Evaluación y pruebas de seguridad
  \item Operaciones de seguridad
  \item Seguridad del desarrollo de software
\end{enumerate}
Casi todos los temas de esta clase (gobernanza, gestión de riesgos, IAM/PAM, evaluación de
seguridad, desarrollo seguro) caen dentro de alguno de estos dominios.
\end{practbox}

\subsection{Vocabulario de seguridad (definiciones de la práctica)}

La práctica fija un vocabulario preciso con el clásico diagrama de Bishop
(\emph{threat agent $\to$ threat $\to$ vulnerability $\to$ risk $\to$ asset}). Conviene
tenerlo claro porque toda la gestión de riesgos se apoya en él.

\begin{center}
\begin{tikzpicture}[font=\footnotesize,
   nodo/.style={draw,black!60,fill=slideGray,rounded corners=1pt,minimum width=2.2cm,
                minimum height=0.7cm,align=center},
   >=Stealth]
  \node[nodo] (ta) {Threat\\agent};
  \node[nodo,right=1.6cm of ta] (th) {Threat\\(amenaza)};
  \node[nodo,right=1.6cm of th] (vu) {Vulnerability\\(vulnerab.)};
  \node[nodo,below=0.9cm of vu] (ri) {Risk\\(riesgo)};
  \node[nodo,left=1.6cm of ri] (as) {Asset\\(activo)};
  \node[nodo,left=1.6cm of as] (ex) {Exposure\\(exposición)};
  \node[nodo,below=0.9cm of ex] (sg) {Safeguard\\(medida)};
  \draw[->] (ta) -- node[above]{da lugar a} (th);
  \draw[->] (th) -- node[above]{explota} (vu);
  \draw[->] (vu) -- node[right]{lleva a} (ri);
  \draw[->] (ri) -- node[above]{daña} (as);
  \draw[->] (as) -- node[above]{causa} (ex);
  \draw[->] (sg) -- node[left]{contrarresta} (ex);
  \draw[->] (sg) |- node[below]{reduce} (ta);
\end{tikzpicture}
\end{center}

\begin{defbox}{Definiciones clave}
\begin{description}[leftmargin=0pt,labelsep=0.4em,style=nextline,itemsep=2pt]
  \item[Amenaza (threat)] Peligro potencial asociado a la explotación de una vulnerabilidad.
  \item[Vulnerabilidad] Falta o debilidad en una medida implementada.
  \item[Riesgo (risk)] \emph{Probabilidad} de que un agente explote la vulnerabilidad y
        produzca efectos en el negocio.
  \item[Exposición] Exposición a pérdidas.
  \item[Medida / control (safeguard)] Mecanismo para reducir el riesgo.
\end{description}
\end{defbox}

\begin{practbox}{Tipos de control (aporte de la práctica)}
Los controles de seguridad se clasifican en tres familias, y la \textbf{defensa en profundidad}
consiste justamente en combinarlas en capas concéntricas:
\begin{itemize}[leftmargin=1.4em,itemsep=1pt]
  \item \textbf{Administrativos} (políticas, procedimientos, capacitación).
  \item \textbf{Técnicos (lógicos)} (firewalls, cifrado, autenticación).
  \item \textbf{Físicos} (puertas, cámaras, control de acceso a instalaciones).
\end{itemize}
\end{practbox}

\begin{practbox}{Políticas vs.\ Mecanismos (aporte de la práctica)}
Distinción fundamental que recorre toda la materia:
\begin{itemize}[leftmargin=1.4em,itemsep=2pt]
  \item \textbf{Política}: define \emph{qué acciones son seguras (permitidas)} y \emph{qué
        acciones son inseguras (prohibidas)}.
  \item \textbf{Mecanismo}: procedimiento, herramienta o método para \emph{garantizar el
        cumplimiento} de la política. Los mecanismos pueden ser de tipo:
        \textbf{prevenir, disuadir, detectar, corregir, recuperar, compensar}.
\end{itemize}
\end{practbox}

%=====================================================================
\section{Gobernanza (Governance)}
%=====================================================================

\begin{defbox}{Gobernanza de seguridad}
Conjunto de \textbf{procesos, políticas, procedimientos y controles} implementados para
\textbf{gestionar y mitigar los riesgos} de la organización, estableciendo
\textbf{responsabilidades y roles} dentro de la entidad para garantizar la protección de la
información.
\end{defbox}

La analogía del profesor: la gobernanza dentro de una empresa es como el \textbf{gobierno} de
un Estado.

\begin{anecbox}{La analogía del gobierno}
\prof{El gobierno expide leyes y después es responsable de alguna forma de hacer que esas
leyes se cumplan. Dentro de una empresa existe algo bastante análogo: la gobernanza está dada
generalmente dentro de un equipo de seguridad, que emite las reglas, decide qué reglas hay que
tomar y se hace responsable de que se cumplan.}
\end{anecbox}

\subsection{Importancia en la organización}
\begin{itemize}[leftmargin=1.4em,itemsep=1pt]
  \item Asegura la confidencialidad, integridad y disponibilidad de la información.
  \item Protege contra amenazas cibernéticas y violaciones de seguridad.
  \item Fomenta la \textbf{confianza de clientes y stakeholders}.
\end{itemize}

\subsection{Objetivos principales de la gobernanza}

\begin{center}
\begin{tikzpicture}[font=\footnotesize]
  \node[circle,fill=itbaCyanDark,text=white,align=center,minimum size=2.0cm,font=\bfseries\scriptsize]
       (c) {CYBER\\SECURITY\\STRATEGY};
  \node[draw=itbaCyanDark,rounded corners,fill=itbaBlue!25,align=center,above left=0.8cm and 1.2cm of c,text width=3.4cm]
       (a) {\textbf{Alinear} la estrategia de seguridad con los objetivos del \textbf{negocio}};
  \node[draw=itbaCyanDark,rounded corners,fill=itbaBlue!25,align=center,above right=0.8cm and 1.2cm of c,text width=3.4cm]
       (b) {\textbf{Reducir} riesgos y \textbf{mitigar} el impacto de incidentes};
  \node[draw=itbaCyanDark,rounded corners,fill=itbaBlue!25,align=center,below left=0.8cm and 1.2cm of c,text width=3.4cm]
       (d) {\textbf{Cumplir} regulaciones y normativas aplicables};
  \node[draw=itbaCyanDark,rounded corners,fill=itbaBlue!25,align=center,below right=0.8cm and 1.2cm of c,text width=3.4cm]
       (e) {Crear una \textbf{cultura de seguridad} en la organización};
  \foreach \n in {a,b,d,e}{\draw[-{Stealth},itbaCyanDark] (c) -- (\n);}
\end{tikzpicture}
\end{center}

\begin{warnbox}{Cuidado: no atentar contra el negocio (aceptación psicológica)}
La estrategia de seguridad \textbf{no debe ir en contra de cómo la empresa gana dinero}. Si se
ponen demasiadas trabas, los empleados \emph{buscan otros caminos} y terminan atentando contra
la seguridad. Es el problema de la \textbf{aceptación psicológica} visto en clases anteriores.
\prof{Si yo le pongo muchas barreras a un empleado [\dots] se van a bajar todos los datos en un
CSV, lo van a publicar en Drive o en algún lugar, y van a empezar a elaborar ahí. Y esto lo
digo porque sé que ha pasado.}
\end{warnbox}

\begin{anecbox}{Por qué importa la cultura: el eslabón más débil}
\prof{¿Cuál es el eslabón más débil en un sistema? Los usuarios, la gente que interactúa con
esos sistemas terminan siendo los más expuestos.} Por eso la gobernanza debe generar una
\textbf{cultura de seguridad}: campañas de concientización, campañas de phishing simuladas y
capacitación, sobre todo a quienes trabajan con tarjetas de crédito, datos personales o
información crítica.
\end{anecbox}

\noindent\textbf{Regulaciones que condicionan la estrategia (ejemplos del profe):} empresas que
cotizan en bolsa deben \emph{informar todo incidente de seguridad} (caso contrario hay
consecuencias penales); el manejo de \textbf{datos personales} exige protegerlos o se arriesgan
\emph{multas millonarias} ante un robo.

%=====================================================================
\section{Gestión de Riesgos (Risk Management)}
%=====================================================================

\begin{defbox}{Gestión de riesgos}
\textbf{Identificación, evaluación y control} de amenazas potenciales para el negocio, con el
objetivo de \textbf{minimizar el impacto} de los incidentes de seguridad en la operación.
\end{defbox}

\begin{anecbox}{Por qué la seguridad se gestiona como un riesgo}
\prof{Antes —y todavía se sigue haciendo— la seguridad se gestionaba como un riesgo más. Las
empresas ya saben de riesgos, saben cómo manejarlos, mitigarlos, eliminarlos. Entonces la
seguridad se trataba como algo más, y la verdad que funciona, sobre todo en las grandes
corporaciones.}
\end{anecbox}

La gestión de riesgos es un \textbf{ciclo continuo} de 4 etapas — nunca termina, porque los
riesgos evolucionan, decrecen, se transforman o aparecen nuevos.

\begin{center}
\begin{tikzpicture}[font=\small\bfseries,>=Stealth]
  \def\R{2.2}
  \node[circle,minimum size=1.5cm,fill=blue!45!cyan,text=white,align=center] (id) at (90:\R) {Identify\\(Identificar)};
  \node[circle,minimum size=1.5cm,fill=red!70,text=white,align=center] (as) at (0:\R) {Assess\\(Evaluar)};
  \node[circle,minimum size=1.5cm,fill=orange!85,text=white,align=center] (mi) at (-90:\R) {Mitigate\\(Mitigar)};
  \node[circle,minimum size=1.5cm,fill=green!55!black,text=white,align=center] (mo) at (180:\R) {Monitor\\(Monitorear)};
  \draw[->,line width=2pt,itbaCyanDark] (id) to[bend left=20] (as);
  \draw[->,line width=2pt,itbaCyanDark] (as) to[bend left=20] (mi);
  \draw[->,line width=2pt,itbaCyanDark] (mi) to[bend left=20] (mo);
  \draw[->,line width=2pt,itbaCyanDark] (mo) to[bend left=20] (id);
\end{tikzpicture}
\end{center}

\subsection{Identificación del riesgo}

Primero se identifican los \textbf{activos críticos para el negocio} (¿qué puede afectar al
ciclo de negocio?) y luego los \textbf{riesgos potenciales} sobre ellos.

\begin{center}
\begin{tabular}{@{}p{0.46\textwidth}p{0.46\textwidth}@{}}
\toprule
\textbf{Activos críticos} & \textbf{Riesgos potenciales} \\
\midrule
\textbullet\ Datos (empleados, clientes, comerciales, financieros, proveedores) &
\textbullet\ Malware \\
\textbullet\ Sistemas (principal activo de las empresas tech) &
\textbullet\ Ataques externos \\
\textbullet\ Infraestructura (servidores, data centers, oficinas físicas) &
\textbullet\ Errores humanos \\
\textbullet\ Personas (las empresas no se mueven solas) &
\textbullet\ Sabotaje (amenaza interna) \\
 & \textbullet\ Ambientales (huracán, incendio) \\
\bottomrule
\end{tabular}
\end{center}

\begin{anecbox}{Riesgos externos a la empresa: caída de una región cloud}
Además de los de la slide, el profe remarca los \textbf{problemas externos}: proveedores de
infraestructura cloud, Internet, DNS\dots
\prof{Se caía una región de AWS y casi todas las empresas argentinas o de Latam están subidas a
esa región [\dots] se caía PedidosYa, no se podía usar.} Si la empresa sólo tiene contratada esa
región, deja de funcionar.
\end{anecbox}

\subsection{Evaluación del riesgo (Assess)}

\begin{itemize}[leftmargin=1.4em,itemsep=1pt]
  \item Se identifica la \textbf{probabilidad de ocurrencia} y el \textbf{impacto} en el negocio.
  \item Se realizan \textbf{evaluaciones externas de seguridad} (Red Team / PenTest).
  \item Se clasifican los riesgos según su impacto: \textbf{Alto / Medio / Bajo}.
\end{itemize}

\begin{anecbox}{La matriz probabilidad $\times$ impacto}
En la práctica se arman \textbf{matrices de probabilidad vs.\ impacto}: se ataca primero lo de
\emph{máxima probabilidad y mayor impacto}, y se barre en diagonal hacia abajo. Hay casos
opuestos: \prof{Es muy poco probable que pase esto, pero si llega a pasar estamos al horno.}
Ahí la empresa, como ente que afronta una pérdida económica, decide qué hacer con cada caso.
\end{anecbox}

\subsection{Mitigación del riesgo (Mitigate)}

En la etapa de mitigación se elaboran \textbf{políticas de gestión}, se definen y evalúan
\textbf{procedimientos de operación y acción ante incidentes}, se diseñan y prueban
\textbf{planes de acción} y se capacita al personal a todo nivel. Existen
\textbf{5 estrategias} para tratar un riesgo:

\begin{center}
\begin{tikzpicture}[font=\footnotesize\bfseries,>=Stealth,
   risk/.style={circle,fill=red!75,text=white,minimum size=1.0cm,font=\scriptsize\bfseries}]
  % ACEPTARLO
  \node[risk] (r1) {RISK};
  \draw[->,black!70,line width=1.2pt] ($(r1.south)+(0,-0.6)$) -- (r1.south);
  \node[below=0.65cm of r1,align=center] {ACEPTARLO};
  % EVITARLO
  \node[risk,right=1.5cm of r1] (r2) {RISK};
  \draw[->,black!70,line width=1.2pt] (r2.south) arc (-90:90:0.55);
  \node[below=0.65cm of r2,align=center] {EVITARLO};
  % TRANSFERIRLO
  \node[risk,right=1.5cm of r2] (r3) {RISK};
  \draw[->,black!70,line width=1.2pt] (r3) ++(0:0.75) arc (0:330:0.75);
  \node[below=0.65cm of r3,align=center] {TRANSFERIRLO};
  % REDUCIRLO
  \node[risk,right=1.5cm of r3] (r4) {RISK};
  \draw[->,black!70,line width=1.2pt] ($(r4.east)+(0.55,0)$) -- (r4.east);
  \draw[->,black!70,line width=1.2pt] ($(r4.west)+(-0.55,0)$) -- (r4.west);
  \node[below=0.65cm of r4,align=center] {REDUCIRLO};
  % COMPENSARLO
  \node[risk,right=1.5cm of r4] (r5) {RISK};
  \draw[->,black!70,line width=1.2pt] ($(r5.north)+(0,0.55)$) -- (r5.north);
  \draw[->,black!70,line width=1.2pt] ($(r5.south)+(0,-0.55)$) -- (r5.south);
  \node[below=0.65cm of r5,align=center] {COMPENSARLO};
\end{tikzpicture}
\end{center}

\begin{defbox}{Las 5 estrategias de tratamiento del riesgo}
\begin{description}[leftmargin=0pt,labelsep=0.4em,style=nextline,itemsep=3pt]
  \item[Evitarlo (avoid)] Quitar la causa del riesgo. Suele ser \emph{costoso y complejo}
        (ej.: arreglar vulnerabilidades en un sistema legacy o migrarlo, con posibles
        \emph{side effects} que rompen el sistema). Se paga el costo si realmente hay que
        solucionarlo.
  \item[Reducirlo (reduce/mitigate)] Bajar el impacto o la probabilidad. \textbf{Es lo más
        común}. Ejemplo: autenticación débil con sólo password $\to$ se le suma un
        \textbf{factor extra (MFA)}; no es tan caro como migrar o dar de baja un sistema.
  \item[Transferirlo (transfer)] Delegarlo a un \textbf{proveedor}; si hay un problema,
        responde el proveedor. Ejemplo: poner Cloudflare como \emph{proxy reverso} y punto de
        entrada de una web expuesta a Internet, aplicando políticas anti-ataque.
  \item[Compensarlo (compensate)] \emph{No es lo mismo que reducir.} Cuando no se puede evitar
        ni reducir, se \textbf{contrarresta el impacto} con un control adicional. Ejemplo:
        dependencias que podrían tener vulnerabilidades y que no se pueden quitar $\to$
        \textbf{auditar y controlar} esa dependencia para detectar CVEs apenas aparecen.
  \item[Aceptarlo (accept)] Asumir el riesgo (típico en riesgos de \emph{baja probabilidad}).
        A veces tomar otra alternativa es más costoso que aceptarlo. \textbf{Siempre se
        documenta}: el equipo de Risk Management deja constancia de que la empresa decidió
        aceptarlo.
\end{description}
\end{defbox}

\begin{anecbox}{Riesgo aceptado en la práctica: el DRP que no se ejecuta}
A las empresas que cotizan en bolsa se les exige tener un \textbf{DRP} (\emph{Disaster Recovery
Plan}). Para los bancos digitales y fintech, el \textbf{Banco Central} exige que el DRP esté
\emph{documentado y probado}. Pero a muchas otras empresas sólo se les exige \emph{tenerlo}:
\prof{Tengo el plan, está aceptado el riesgo, pero no lo ejecuto porque ejecutarlo me sale un
montón de plata. Si pasa algo, pasa, y confiamos en que funcione.} Ejecutar un DRP implica tener
toda la infra duplicada con todos los enlaces y aplicaciones levantadas, lo cual es muy costoso.
\end{anecbox}

\subsection{Monitoreo y Evaluación (Monitor)}

\begin{itemize}[leftmargin=1.4em,itemsep=1pt]
  \item \textbf{Elaboración y supervisión de controles y medidas de seguridad}.
  \item Auditorías internas y externas.
  \item Simulaciones de ataques e incidentes.
\end{itemize}

Es la etapa que \textbf{cierra el ciclo}: como la organización es un \emph{sistema vivo}, hay
que estar ejecutando y revisando esto constantemente. Existen \textbf{frameworks} que dan un
punto de partida; el profe menciona \textbf{NIST} (la organización de estándares de EE.UU.,
análoga al IRAM en Argentina), que desarrolló frameworks para medir, mitigar y gestionar
riesgos de forma segura.

\begin{slidebox}
\textbf{CIS Critical Security Controls v8} (\url{https://www.cisecurity.org/controls}) — la
slide muestra los 18 controles agrupados; los \emph{foundational} (1--8) son: \emph{(1)}
Inventory \& Control of Assets, \emph{(2)} Inventory \& Control of Software, \emph{(3)} Data
Protection, \emph{(4)} Secure Configuration, \emph{(5)} Account Management, \emph{(6)} Access
Control, \emph{(7)} Continuous Vulnerability Management, \emph{(8)} SIEM. Los restantes incluyen
\emph{(14)} Security Awareness \& Training, \emph{(16)} Application Software Security,
\emph{(17)} Incident Response, \emph{(18)} Penetration Testing.
\end{slidebox}

\begin{slidebox}
\textbf{Conclusión (slide):} la gestión de riesgos es \textbf{fundamental} para proteger la
información; un \textbf{enfoque proactivo} reduce la probabilidad y el impacto de incidentes; y
la \textbf{revisión y adaptación continua} es esencial para enfrentar nuevos desafíos.
\end{slidebox}

%=====================================================================
\section{Políticas de Seguridad}
%=====================================================================

\begin{anecbox}{La política como ``Biblia'' de seguridad}
\prof{La política de seguridad es como una especie de Biblia para la seguridad de la empresa.
Es un documento bastante extenso y denso que se confecciona una sola vez cada cierto tiempo,
se va revisando y lo aprueba el director de seguridad (CISO).} Es muy genérica y abarcativa;
por eso luego se \textbf{desglosa} en documentos cada vez más específicos hasta llegar al
empleado más raso.
\end{anecbox}

La jerarquía de documentos va de \emph{mayor ownership del management / menor especificidad}
(arriba) a \emph{menor ownership / mayor especificidad} (abajo). Las dos capas superiores son
\textbf{mandatorias}; las inferiores, \textbf{discrecionales}.

\begin{center}
\begin{tikzpicture}[font=\footnotesize,>=Stealth]
  % pirámide de 5 niveles
  \fill[red!55!black]   (0,4)   -- (-1.0,3) -- (1.0,3) -- cycle;
  \fill[red!70]         (-1.0,3) -- (1.0,3) -- (1.8,2.2) -- (-1.8,2.2) -- cycle;
  \fill[red!85]         (-1.8,2.2) -- (1.8,2.2) -- (2.6,1.4) -- (-2.6,1.4) -- cycle;
  \fill[orange!80]      (-2.6,1.4) -- (2.6,1.4) -- (3.4,0.6) -- (-3.4,0.6) -- cycle;
  \fill[yellow!75!orange](-3.4,0.6) -- (3.4,0.6) -- (4.2,-0.2) -- (-4.2,-0.2) -- cycle;
  \node[white,font=\bfseries\scriptsize,align=center] at (0,3.45) {POLICY};
  \node[white,font=\bfseries\scriptsize,align=center] at (0,2.6) {STANDARDS};
  \node[white,font=\bfseries\scriptsize,align=center] at (0,1.8) {PROCEDURES};
  \node[white,font=\bfseries\scriptsize,align=center] at (0,1.0) {GUIDELINES};
  \node[black,font=\bfseries\scriptsize,align=center] at (0,0.2) {BASELINES};
  % flechas laterales
  \draw[<->,itbaCyanDark,line width=1pt] (-4.8,-0.2) -- (-4.8,4)
       node[midway,left,align=center,rotate=90,yshift=8pt]{\scriptsize Ownership del management $+$};
  \draw[<->,itbaCyanDark,line width=1pt] (4.8,-0.2) -- (4.8,4)
       node[midway,right,align=center,rotate=-90,yshift=8pt]{\scriptsize Especificidad $+$ (hacia abajo)};
\end{tikzpicture}
\end{center}

\begin{defbox}{Niveles de la jerarquía de políticas}
\begin{description}[leftmargin=0pt,labelsep=0.4em,style=nextline,itemsep=2pt]
  \item[Policy (Política)] \emph{General management statement}. Documento genérico y
        abarcativo de toda la seguridad de la organización. \textbf{Mandatorio}.
  \item[Standards (Estándares)] \emph{Specific mandatory controls}. Reglas obligatorias de
        seguridad que el hardware y el software deben cumplir. \textbf{Mandatorio}. Ej.: un
        estándar de desarrollo seguro.
  \item[Procedures (Procedimientos)] \emph{Step-by-step instructions}. Tareas detalladas paso
        a paso para lograr algún objetivo.
  \item[Guidelines (Guías / directrices)] \emph{Recommendations / best practices}. Acciones
        recomendadas para usuarios y personal cuando no puede aplicarse un estándar específico.
        \textbf{Discrecional}.
  \item[Baselines] \emph{Uniform ways for safeguard implementation}. Suelen tomar forma de
        cursos, capacitaciones y concientizaciones. \textbf{Discrecional}.
\end{description}
\end{defbox}

\begin{practbox}{Estándares de NIST (aporte de la práctica)}
La práctica detalla los estándares concretos del \textbf{NIST} (\emph{National Institute of
Standards and Technology}):
\begin{itemize}[leftmargin=1.4em,itemsep=2pt]
  \item \textbf{NIST 800-53}: \emph{obligatorio} para sistemas de información federales (EE.UU.).
  \item \textbf{NIST CSF 2.0} (\emph{Cybersecurity Framework}): \emph{voluntario} para las
        organizaciones, basado en riesgos, más flexible y adaptable. Sus funciones son:
        \textbf{Gobernar, Identificar, Proteger, Detectar, Responder, Recuperar}. Distingue
        \emph{Gobernanza} (cómo se gestiona la organización) de \emph{Gobernanza de seguridad}
        (cómo se gestiona la seguridad vía políticas y procedimientos).
  \item \textbf{NIST SP 800-207}: define la \textbf{Zero Trust Architecture (ZTA)}, con las
        ideas \emph{Never Trust, Always Verify, Assume breach}. (Lectura recomendada de la
        clase.)
\end{itemize}
\end{practbox}

%=====================================================================
\section{Estrategias de seguridad}
%=====================================================================

Con el tiempo surgieron formas genéricas de \emph{implementar} las políticas de seguridad. Como
las políticas de empresas similares se parecen mucho entre sí (\prof{la política de Mercado
Libre debe ser muy parecida a la de OLX o PedidosYa}), aparecieron estrategias reutilizables.
Hoy se destacan \textbf{dos}: \emph{Defense-in-depth} y \emph{Zero Trust Architecture}.

\subsection{Defense-in-depth (Defensa en profundidad)}

\begin{defbox}{Defense-in-depth}
Estrategia que aplica \textbf{múltiples capas de defensa} para proteger los activos de la
organización. Parte de asumir que \textbf{ningún control de seguridad es perfecto} (la seguridad
nunca está asegurada al 100\%) y busca crear \textbf{redundancia}: cada capa está pensada para
\emph{contener a la capa que sigue}.
\end{defbox}

\begin{anecbox}{La cebolla}
\prof{Esto se piensa como si fuera una cebolla: agarrando capas y metiéndote cada vez más
adentro hasta llegar a lo más importante. Si falla la capa de arriba, la contiene la de abajo,
y así.}
\end{anecbox}

Las \textbf{5 capas} (de afuera hacia el activo crítico):

\begin{center}
\begin{tikzpicture}[font=\footnotesize\bfseries]
  \def\cx{0}\def\cy{0}
  \fill[purpleZTA!85] (\cx,\cy) circle (3.2);
  \fill[purpleZTA!65] (\cx,\cy-0.35) circle (2.6);
  \fill[purpleZTA!45!itbaCyan] (\cx,\cy-0.7) circle (2.0);
  \fill[itbaCyan!75] (\cx,\cy-1.05) circle (1.4);
  \fill[itbaCyanDark] (\cx,\cy-1.4) circle (0.8);
  \node[white] at (\cx,\cy+2.7) {Perimeter};
  \node[white] at (\cx,\cy+1.65) {Network};
  \node[white] at (\cx,\cy+0.6) {Host};
  \node[white] at (\cx,\cy-0.5) {Application};
  \node[white] at (\cx,\cy-1.4) {Data};
\end{tikzpicture}
\end{center}

\begin{defbox}{Las 5 capas de Defense-in-depth}
\begin{description}[leftmargin=0pt,labelsep=0.4em,style=nextline,itemsep=2pt]
  \item[Perímetro] Ej.: un \textbf{firewall} que deniega llamadas de Internet y sólo acepta de
        cierta red o VPN.
  \item[Red (Network)] \textbf{Segmentación de red}: separar la red corporativa (operaciones)
        de la red de servidores. Un empleado del call center no debe poder llegar al servidor
        de ventas.
  \item[Host] \textbf{``Jardinizar''} (hardening) los hosts: el host tiene un software fijo,
        inmutable, no modificable —como un contenedor Docker—, optimizado con buenas prácticas.
  \item[Aplicación] Buen diseño, buenas prácticas, sin vulnerabilidades (todo lo visto en las
        clases anteriores de la materia).
  \item[Datos] La capa más interna, lo que más importa. Para llegar, un atacante debe atravesar
        \emph{todas} las capas anteriores: firewall $\to$ moverse lateralmente entre redes $\to$
        vulnerar el host jardinizado $\to$ vulnerar la aplicación.
\end{description}
\end{defbox}

\paragraph{Soluciones por capa (corte de la slide ``Outside Threat Protection'').}
La diapositiva muestra, en abanico, qué herramientas pueblan cada capa, atravesadas por las
funciones \emph{Prevention, Policy Management, Operations, Monitoring \& Response}:

\begin{center}
\begin{tabular}{@{}>{\bfseries}p{0.22\textwidth}p{0.72\textwidth}@{}}
\toprule
Capa & Herramientas (de la slide y comentadas por el profe) \\
\midrule
Perímetro & \textbf{Honeypot} (deriva al atacante a una ``red boba'' para entretenerlo),
            Message Security (anti-virus/anti-malware), Secure DMZs, Perimeter Firewall,
            Perimeter IDS/IPS. \\
Red & \textbf{NAC} (Network Access Control: decide si un dispositivo se conecta según IP, MAC,
       firmware o credenciales), Web Proxy / Content Filtering (ej.\ Fortinet: bloquea Netflix
       o sitios maliciosos), Enterprise Message Security, \textbf{DLP}. \\
Host (Endpoint) & Content Security (anti-virus/anti-malware), Host IDS/IPS, Endpoint Security
       Enforcement, FDCC Compliance, Patch Management. \\
Aplicación & \textbf{WAF} (Web Application Firewall: como un firewall pero para tráfico
       HTTP/HTTPS de una app web; filtra o reemplaza tráfico según reglas), Dynamic App
       Testing, Database Monitoring/Scanning. \\
Datos & \textbf{IAM} (Identity \& Access Management: autenticación, autorización, gestión de
       identidad, roles, permisos), PKI, Data Classification, Data Wiping, Data/Drive
       Encryption, Database Secure Gateway. \\
Monitoring & SOC/NOC 24x7, CIRT (Incident Reporting/Detection/Response), \textbf{SIEM}, Digital
       Forensics. \\
\bottomrule
\end{tabular}
\end{center}

\begin{defbox}{IDS / IPS y DLP}
\begin{description}[leftmargin=0pt,labelsep=0.4em,style=nextline,itemsep=2pt]
  \item[IDS / IPS] \emph{Intrusion Detection/Prevention System}. Detectan patrones tipo
        antivirus (por \textbf{firmas} o \textbf{anomalías}) y deciden si bloquear, alertar o
        actuar.
  \item[DLP] \emph{Data Loss Prevention}. \prof{No es que se pierda data dentro de la compañía,
        sino que no se filtre data hacia afuera.} Escanean para que no se exfiltren datos de
        tarjeta o personales; si pasa, levantan alertas.
\end{description}
\end{defbox}

\subsubsection{Principios de Defense-in-depth}
\begin{itemize}[leftmargin=1.4em,itemsep=2pt]
  \item \textbf{Diversificación}: usar diversas tecnologías y métodos $\Rightarrow$ una sola
        vulnerabilidad no atraviesa todo.
  \item \textbf{Redundancia}: múltiples capas que se superponen para asegurar protección continua.
  \item \textbf{Segmentación}: dividir la red en segmentos controlados para limitar el acceso no
        autorizado. \prof{Hay compañías muy avanzadas donde tu aplicación no ve nada: si sólo
        tiene que hablar con otra aplicación, sólo tiene visibilidad de esa otra.}
  \item \textbf{Monitoreo}: visibilidad de todas las capas y \textbf{correlación de eventos}
        para detectar ataques.
\end{itemize}

\begin{warnbox}{Problemas de Defense-in-depth}
\begin{itemize}[leftmargin=1.4em,itemsep=1pt]
  \item \textbf{Altos costos}: de tecnología (muchos proveedores distintos, $\geq 5$
        presupuestos), de capacitación y operación (no hay ``bala de plata'', cada producto es
        distinto), y de monitoreo (\emph{se recolectan muchos datos pero se les saca poco
        valor}).
  \item \textbf{No resuelve eficientemente los ataques internos.} Si alguien ya está adentro
        (robó credenciales de un empleado), es muy difícil darse cuenta: no todas las
        herramientas son bidireccionales.
  \item \textbf{Difícil de implementar con sistemas distribuidos}: data centers híbridos
        (propios + terceros), múltiples proveedores de nube, empresas \emph{cloud native}. Esta
        estrategia fue pensada para la época de \emph{data center / on-premise}.
\end{itemize}
\end{warnbox}

\begin{anecbox}{El intranet ``de confianza'' que quedó obsoleto}
Tradicionalmente, dentro del \textbf{intranet} de la empresa se relajaban los controles (menos
autenticación, sesiones más largas) para no friccionar al empleado. \prof{Eso durante mucho
tiempo se consideró buena práctica y seguro. Hoy sabemos que estamos mirando otro tipo de
estrategias.} Eso da pie al Zero Trust.
\end{anecbox}

\subsection{Zero Trust Architecture (ZTA)}

\begin{defbox}{Zero Trust Architecture (NIST SP 800-207)}
Modelo de seguridad que asume que las amenazas, \textbf{tanto internas como externas}, están
presentes en la red \textbf{en todo momento}. \textbf{Objetivo:} proteger la integridad,
confidencialidad y disponibilidad de los datos y recursos. \textbf{Premisa:} cualquier red
interna es tan peligrosa como Internet (``Internet es \emph{Bad Lands}, la tierra de nadie'').
\end{defbox}

\begin{anecbox}{Por qué ninguna red es segura}
\prof{Si vos fuiste a Starbucks, te fuiste al baño, dejaste la máquina y vino un chabón con un
pendrive y te instaló un exploit/malware, y venís y te conectás\dots yo no te pido nada, asumo
que tu computadora está bien y ya estás inyectando algo dentro de la red.} Con el trabajo remoto
es aún peor: no se sabe quién accede. Por eso se considera que \textbf{toda red interna es tan
peligrosa como Internet} y se aplican adentro los mismos controles que afuera.
\end{anecbox}

\begin{defbox}{Los 3 principios básicos de Zero Trust}
\begin{enumerate}[leftmargin=1.6em,itemsep=1pt]
  \item \textbf{Nunca confíes} (\emph{Never Trust}).
  \item \textbf{Siempre verifica} (\emph{Always Verify}).
  \item \textbf{Asume que has sido atacado} (\emph{Assume breach}).
\end{enumerate}
\end{defbox}

\begin{anecbox}{Cómo se relaja en la práctica: SSO / federación}
En teoría \emph{cada} componente debería forzar la verificación, pero pedir 20 logins a 20
aplicaciones haría que el empleado se queje. En la práctica se relaja con \textbf{federación de
identidad}: \prof{te federo el login de una sola forma, por ejemplo a través de Active
Directory, lo hacés una vez; después cada aplicación tiene los logins federados, no te pide un
login nuevo, pero los controles de seguridad se aplican porque se heredan los permisos y roles
de esa persona en cada pieza.}
\end{anecbox}

\subsubsection{Componentes de ZTA (modelo NIST 800-207)}

El núcleo separa un \textbf{Control Plane} (donde vive el \textbf{PDP} = Policy Decision Point,
formado por el \emph{Policy Engine} y el \emph{Policy Administrator}) de un \textbf{Data Plane}
(donde está el \textbf{PEP} = Policy Enforcement Point). El sujeto accede a través de un sistema
considerado \emph{untrusted}; sólo tras pasar el PEP el acceso al recurso es \emph{trusted}.

\begin{center}
\begin{tikzpicture}[font=\scriptsize,>=Stealth,
   blk/.style={draw=black!50,fill=slideGray,rounded corners=1pt,minimum height=0.65cm,minimum width=1.9cm,align=center},
   red/.style={draw=attackRed,fill=attackRed!75,text=white,rounded corners=2pt,minimum height=0.75cm,minimum width=2.0cm,align=center,font=\bfseries\scriptsize}]
  % columna izquierda (entradas) — arriba
  \node[blk] (cdm) at (0,3.6) {CDM System};
  \node[blk] (ic)  at (0,2.7) {Industry\\Compliance};
  \node[blk] (ti)  at (0,1.8) {Threat\\Intelligence};
  \node[blk] (al)  at (0,0.9) {Activity Logs};
  % columna derecha (entradas) — arriba
  \node[blk] (dap)  at (10,3.6) {Data Access\\Policy};
  \node[blk] (pki)  at (10,2.7) {PKI};
  \node[blk] (idm)  at (10,1.8) {ID Management};
  \node[blk] (siem) at (10,0.9) {SIEM System};
  % control plane (centro arriba)
  \node[red] (pe) at (5,3.1) {Policy Engine};
  \node[red] (pa) at (5,2.0) {Policy\\Administrator};
  \node[draw=black!40,dashed,fit=(pe)(pa),inner sep=6pt,label=above:{\bfseries Control Plane (PDP)}] (cp) {};
  % data plane (fila inferior, ancho completo)
  \node[blk,minimum width=1.5cm] (subj) at (0.4,-0.7) {Subject};
  \node[blk,minimum width=1.5cm] (sys)  at (2.6,-0.7) {System};
  \node[red,minimum width=2.4cm] (pep)  at (5.6,-0.7) {Policy\\Enforcement Point};
  \node[draw=black!50,fill=itbaBlue!30,rounded corners,minimum height=0.7cm,align=center] (er) at (9.0,-0.7) {Enterprise\\Resource};
  \node[font=\scriptsize\itshape] at (5.6,-1.7) {\textbf{Data Plane}};
  % flechas de las columnas hacia el centro
  \foreach \n in {cdm,ic,ti,al}{\draw[->,black!55] (\n.east) -- ++(0.5,0);}
  \foreach \n in {dap,pki,idm,siem}{\draw[<-,black!55] (\n.west) -- ++(-0.5,0);}
  % flujo data plane
  \draw[->,black!55] (subj) -- (sys);
  \draw[->,black!55] (sys) -- node[above]{Untrusted} (pep);
  \draw[->,black!55] (pep) -- node[above]{Trusted} (er);
  % control -> enforcement
  \draw[->,dashed,black!55] (pa.south) -- (pep.north);
  \draw[->,dashed,black!55] (pe) -- (pa);
\end{tikzpicture}
\end{center}

\begin{defbox}{Componentes que alimentan al motor de políticas}
\begin{description}[leftmargin=0pt,labelsep=0.4em,style=nextline,itemsep=2pt]
  \item[CDM — Continuous Diagnostics and Mitigation] Sistema de \emph{diagnóstico y mitigación
        continuos}. Provee al motor de políticas información sobre el activo que solicita acceso:
        ¿tiene el antivirus actualizado? ¿los últimos updates de seguridad? Si una regla no se
        cumple, puede \textbf{desactivar el recurso} (ej.: software en la PC del empleado que no
        se puede desinstalar y, ante una anomalía, deshabilita la máquina).
  \item[Industry Compliance] Garantiza que la empresa cumpla los regímenes normativos
        aplicables: \textbf{HIPAA} (datos médicos), \textbf{Banco Central}, \textbf{SOX}
        (cotizar en bolsa, transparencia), \textbf{PCI} (datos de tarjeta de crédito).
  \item[Threat Intelligence] Fuentes de inteligencia sobre amenazas (internas/externas) que
        informan sobre nuevos ataques y vulnerabilidades (CVEs, hashes de malware, IPs
        comprometidas, botnets). Ej.\ de la slide: \emph{VirusTotal}, feeds de IoC.
  \item[Activity Logs] Registros de actividad de la red y el sistema: tráfico, accesos a
        recursos y eventos que dan una visión (casi) en tiempo real de la postura de seguridad.
  \item[SIEM — Security Information and Event Management] Recolecta eventos y los
        \textbf{correlaciona} para su posterior análisis, detectar anomalías y generar alertas.
  \item[ID Management] Crea, almacena y gestiona las cuentas de usuario y los registros de
        identidad (nombre, email, certificados, función, privilegios y derechos). \prof{Yo soy
        Nando Suárez, pero en MELI soy lee.suarez: una persona puede tener varios usuarios.}
  \item[PKI — Public Key Infrastructure] Genera y registra los certificados emitidos por la
        empresa a recursos, sujetos, servicios y aplicaciones.
  \item[Data Access Policy] Características, reglas y políticas sobre el acceso a los recursos.
        Pueden estar codificadas (vía interfaz de gestión) o generarse dinámicamente por el PE,
        y son el punto de partida para autorizar el acceso a un recurso.
\end{description}
\end{defbox}

\begin{slidebox}
\textbf{ZTA en el mundo real:} todos los grandes proveedores publican guías de implementación —
Azure, AWS, RedHat (Linux), Windows Server. La adopción de ZTA está en plena ``oleada''.
\end{slidebox}

\begin{warnbox}{Cuidado: ZTA mal implementado}
ZTA parece la panacea, pero \prof{es muy fácil que lo implementemos mal.} Casos típicos:
implementarlo \emph{de forma deficiente} (mal informados), o implementarlo y \textbf{no
controlar nada} —``tenemos ZTA'' pero no se monitorea ni se hace seguimiento de cómo evoluciona.
Por eso son importantes las guías/framework y las \textbf{auditorías}.
\prof{Si arrancás de cero, casi seguro conviene ir por ZTA, porque es arrancar sin ningún sesgo
ni \emph{baggage}.} Migrar un Defense-in-depth existente a ZTA es muy difícil.
\end{warnbox}

%=====================================================================
\section{Privileged Access Management (PAM)}
%=====================================================================

\begin{defbox}{Privileged Access Management (PAM)}
Se utiliza para \textbf{proteger y gestionar el uso de cuentas especiales} de una organización
que tienen \textbf{acceso privilegiado} a sistemas críticos. Estos accesos son especialmente
sensibles porque pueden: conceder a los usuarios capacidades amplias, modificar configuraciones
de seguridad, acceder a datos confidenciales y supervisar otras cuentas de usuario.
\end{defbox}

\begin{anecbox}{Por qué los atacantes buscan estas cuentas}
\prof{Estos usuarios son lo que más buscan los atacantes, porque una vez que te haces con uno de
ellos ya tenés allanado el camino para hacer daño, robar cosas, y tenés menos trabas para
ejecutar tus acciones.}
\end{anecbox}

\subsection{Prácticas y herramientas de PAM}

\begin{enumerate}[leftmargin=1.6em,itemsep=3pt]
  \item \textbf{Inventario de cuentas privilegiadas.} Inventariar todas las cuentas con acceso
        privilegiado para asegurar \emph{visibilidad y control}; documentarlas y registrarlas.
  \item \textbf{Principio de privilegio mínimo.} Los usuarios deben tener el mínimo acceso
        necesario para su trabajo. La cuenta privilegiada \textbf{no es la cuenta de trabajo
        diario}: es una cuenta separada, que puede transferirse cuando la persona se va o cambia
        de rol (de ahí que se reutilicen, lo cual tiene un lado malo).
  \item \textbf{Autenticación multifactor (MFA).} 2, 3 o más factores para garantizar que sólo
        usuarios autorizados accedan, y \textbf{el no repudio} (nadie puede negar haber actuado:
        ``vos tenías el TOTP, sólo vos podías avanzar sobre esos factores'').
  \item \textbf{Registro y monitoreo.} Se registran y monitorean todas las actividades —en un
        servidor externo, guardando logs de auditoría— para detectar actividad sospechosa y
        poder \emph{reconstruir lo que pasó}.
  \item \textbf{Rotación de contraseñas.} Las contraseñas se cambian regular y automáticamente
        para reducir el riesgo de \emph{compromisos prolongados}. También se hacen
        \textbf{reválidas}: cuando una cuenta se transfiere se cambia la contraseña, y cada
        tanto se revalida que la persona deba seguir teniéndola.
\end{enumerate}

\subsection{Ciclo de vida de PAM (PAM Life cycle)}

\begin{center}
\begin{tikzpicture}[font=\scriptsize\bfseries,>=Stealth]
  \def\R{2.6}
  \foreach \ang/\lbl/\i in {90/Define/1, 30/Discover/2, -30/{Manage \&\\Protect}/3,
                            -90/Monitor/4, -150/{Detect Abnormal\\Usage}/5, 150/{Respond to\\Incidents}/6}{
    \node[circle,fill=purpleZTA!85,text=white,minimum size=1.25cm,align=center,font=\tiny\bfseries] (n\i) at (\ang:\R) {\lbl};
  }
  \node[circle,draw=practGreen,line width=1.5pt,text=practGreen,minimum size=1.0cm,font=\bfseries] (c) at (0,0) {PAM};
  % Review & Audit (séptimo, arriba-izq entre Respond y Define)
  \foreach \a/\b in {1/2,2/3,3/4,4/5,5/6}{\draw[->,line width=1.2pt,practGreen] (n\a) to[bend left=12] (n\b);}
  \draw[->,line width=1.2pt,practGreen] (n6) to[bend left=12] (n1);
\end{tikzpicture}
\end{center}

\noindent El ciclo (según la slide): \textbf{Define $\to$ Discover $\to$ Manage \& Protect $\to$
Monitor $\to$ Detect Abnormal Usage $\to$ Respond to Incidents $\to$ Review \& Audit} (y vuelve
a empezar). Resumido por el profe como: \emph{descubrir, manejar y proteger, monitorear,
detectar usos anormales, respuesta a incidentes, revisar y auditar}.

%=====================================================================
\section{Prevención de vulnerabilidades: DevSecOps}
%=====================================================================

Muchas organizaciones \textbf{construyen su propio software} y operan a través de él. A
diferencia de los productos que se compran (donde el proveedor vela por las vulnerabilidades),
aquí \emph{nosotros} debemos gestionar y prevenir las vulnerabilidades de nuestro código.

\begin{anecbox}{De DevOps a DevSecOps (la historia que contó el profe)}
\textbf{Antes (80s/90s/2000s):} el equipo de \emph{desarrollo} programaba y un equipo de
\emph{operaciones (Ops)} desplegaba y monitoreaba —aislados y con \textbf{objetivos opuestos}:
desarrollo quería sacar features rápido, operaciones quería no tocar nada para no romper la
estabilidad. \prof{Se bloqueaba todo, había peleas, enfrentamientos.}%
\textbf{DevOps:} alguien dice ``estos equipos ahora son uno solo'', con el mismo objetivo: ser
responsables de todo lo que se despliega y de que funcione. De ahí surgen CI/CD, infraestructura
como código y las \emph{métricas} (New Relic, Datadog, Grafana). Lema: \prof{\emph{you build it,
you run it}}%
\textbf{DevSecOps:} se aplica el mismo principio al \emph{equipo de seguridad}, que también
estaba aislado y quería ``frenar todo''. La idea: meter la seguridad en \textbf{cada paso} del
proceso.
\end{anecbox}

\begin{defbox}{DevSecOps = Development + Security + Operations}
\textbf{Metodología} que integra la seguridad en \emph{todos} los procesos de desarrollo de
software y operaciones de TI. La idea central es el \textbf{Shift Left}: meter la seguridad cada
vez más ``a la izquierda'' del proceso (desde la planificación), porque \textbf{cuanto más
temprano se detecta un problema, menos costoso es resolverlo}.
\end{defbox}

\begin{center}
\begin{tikzpicture}[font=\scriptsize,>=Stealth,
  fase/.style={signal,signal to=east,draw=black!40,minimum height=0.85cm,minimum width=2.3cm,align=center,font=\bfseries\scriptsize}]
  \node[fase,fill=red!25]    (f1) {PLAN \&\\DEVELOP};
  \node[fase,fill=purple!20,right=0.15cm of f1] (f2) {COMMIT\\THE CODE};
  \node[fase,fill=orange!25,right=0.15cm of f2] (f3) {BUILD \&\\TEST};
  \node[fase,fill=blue!18,right=0.15cm of f3]   (f4) {GO TO\\PRODUCTION};
  \node[fase,fill=green!22,right=0.15cm of f4]  (f5) {OPERATE};
  \draw[->,line width=2.5pt,attackRed] ($(f5.south)+(0,-0.45)$) -- node[below,font=\bfseries,attackRed]{Shift Left} ($(f1.south)+(0,-0.45)$);
\end{tikzpicture}
\end{center}

\begin{anecbox}{Security Champions: cómo escalar seguridad en cada equipo}
El equipo de seguridad puede ser ``10 personas vs.\ 1500 desarrolladores''. ¿Cómo estar en cada
equipo? Con el programa \textbf{Security Champions}: se delega en una persona de cada equipo el
rol de ``campeón de seguridad''. No es trivial (hay que capacitar, y el expertise no es el de un
especialista). \prof{DevSecOps implica que haya gente de seguridad en cada uno de los equipos
desde el primer momento en que se toma un desarrollo y se planifica.} Muchas empresas dicen tener
DevSecOps cuando en realidad sólo hacen herramientas que le piden usar a los desarrolladores.
\end{anecbox}

\subsection{Las 5 fases de DevSecOps}

\begin{description}[leftmargin=0pt,labelsep=0.4em,style=nextline,itemsep=4pt]
  \item[1. Plan \& Develop (Planeamiento y Desarrollo)]
    Aplicar todos los principios de seguridad vistos en la materia.
    \begin{itemize}[leftmargin=1.4em,itemsep=0pt,topsep=2pt]
      \item \textbf{Threat modeling}: evaluar las amenazas/riesgos de una solución o feature
            nuevo \emph{antes} de desarrollarlo (nueva superficie de ataque).
      \item \textbf{IDE security plugins} y uso de asistentes de IA para código seguro.
      \item \textbf{Pre-commit hooks}: forzar controles \emph{antes} del commit, corriendo
            localmente (que no se expongan secretos ni vulnerabilidades).
      \item \textbf{Secure coding standards}, capacitación, \textbf{revisión de PRs} (peer
            review) y revisión de vulnerabilidades/correcciones.
    \end{itemize}
  \item[2. Commit the code (On Commit/PR)]
    Controles \textbf{automatizados} antes de que el código siga el deployment.
    \begin{itemize}[leftmargin=1.4em,itemsep=0pt,topsep=2pt]
      \item \textbf{SAST} (Static Application Security Testing): análisis estático del código
            (ej.\ SonarQube), que marca \emph{code smells}, malas prácticas y problemas de
            seguridad.
      \item \textbf{Security unit \& functional tests}: unit tests específicos de seguridad,
            con foco en los \textbf{CWE}; probar inyecciones y vulnerabilidades vistas en clase.
      \item \textbf{Dependency management}: control de vulnerabilidades en las dependencias.
      \item \textbf{Secure pipelines}: proteger el acceso a recursos confidenciales y claves de
            seguridad. Los pipelines definen, como antes la \emph{cobertura de testing}, ahora
            un \emph{nivel de seguridad} del feature.
    \end{itemize}
  \item[3. Build \& Test (On Build/Test, ambientes bajos / sandbox)]
    \begin{itemize}[leftmargin=1.4em,itemsep=0pt,topsep=2pt]
      \item \textbf{DAST} (Dynamic Application Security Testing): controles dinámicos sobre la
            app \emph{ejecutándose}, vía sus interfaces (APIs, URL endpoints).
      \item \textbf{Cloud configuration validation} e \textbf{Infrastructure scanning}: validar
            el código que crea recursos en la nube (IaC) y controlar la configuración tras
            correr el IaC.
      \item \textbf{Security acceptance testing}: evaluación de seguridad de la aplicación como
            un todo; ejecución de ataques (mini pen-tests automatizados).
    \end{itemize}
  \item[4. Go to Production (On Deployment)]
    \begin{itemize}[leftmargin=1.4em,itemsep=0pt,topsep=2pt]
      \item \textbf{Security smoke tests} y \textbf{configuration checks} sobre los artefactos
            consumidos en el despliegue (configuraciones, secretos). Las configuraciones de
            producción difieren de las de ambientes bajos.
      \item Evaluación de la app \emph{incluyendo el ambiente de despliegue} (Firewalls, WAF,
            IPS).
      \item \textbf{Live site penetration testing} (pen-tests externos) y aparición de los
            \textbf{Red / Blue Team}.
    \end{itemize}
  \item[5. Operate (On Operation)]
    \begin{itemize}[leftmargin=1.4em,itemsep=0pt,topsep=2pt]
      \item \textbf{Continuous monitoring} de métricas y logs de seguridad.
      \item Integración con \textbf{Threat Intelligence}.
      \item \textbf{Blind Penetration Test}, ejercicios de simulación de incidentes y
            \textbf{respuesta ante incidentes}.
      \item \textbf{Blameless postmortems}.
    \end{itemize}
\end{description}

\begin{anecbox}{El estado real de la industria}
\prof{Hoy la gran mayoría de las empresas tiene la seguridad atacando esta última etapa
(operación): ``ya se desarrolló, ya se testeó, ya está, ahora quiero ver la seguridad''. La idea
es que la seguridad surja desde el principio. Cuanto más al principio, mejor.}
\end{anecbox}

\subsection{Testeos de seguridad en el código}

\begin{practbox}{DAST / SAST / SCA (la práctica los agrupa explícitamente)}
\begin{description}[leftmargin=0pt,labelsep=0.4em,style=nextline,itemsep=2pt]
  \item[DAST — Dynamic Application Security Testing] Evalúa la seguridad de la app
        \textbf{mientras se ejecuta}, interactuando como un usuario externo. Identifica
        vulnerabilidades explotables en la red (validación de entrada, autenticación). Usa las
        interfaces estándares de la app; se corre en ambientes bajos y \textbf{automatiza
        pruebas repetibles} (escanear puertos, inputs inválidos, requests que generan errores,
        clásico \emph{Stack trace} expuesto).
  \item[SAST — Static Application Security Testing] Evalúa la app \textbf{analizando el código}
        (estático), sin ejecutarla.
  \item[SCA — Software Composition Analysis] Analiza las \textbf{dependencias y bibliotecas}
        usadas para identificar vulnerabilidades de seguridad, problemas de \emph{licencia} y
        otros riesgos (incluidas dependencias \emph{transitivas}). Debe extenderse a los
        servicios SaaS consumidos (bases de datos online, versiones de APIs, ambientes cloud).
\end{description}
\end{practbox}

\begin{anecbox}{Log4Shell: por qué importa el SCA y el SBOM}
\prof{A mí me tocó vivir la de Log4Shell.} Vulnerabilidad crítica (2021) que permitía
\textbf{ejecución remota de código} a través de \emph{Log4j}, una librería de logging usada por
casi cualquier aplicación Java del planeta. \prof{Salían todas las PoC [\dots] cosas
impresionantes, y todos salieron a correr a ver qué tenía su aplicación, porque nadie sabía si
realmente estaba usando Log4j: era una dependencia tan estándar que uno se olvidaba que la
tenía.} El parche salió rápido, pero \emph{todas} las apps del mundo tuvieron que actualizar y
desplegar. Moraleja: sin un inventario (SBOM) no sabés si estás afectado.
\end{anecbox}

\subsection{Inventario de software: SBOM}

\begin{defbox}{SBOM — Software Bill of Materials}
Inventario \textbf{formal y legible por máquina} de los componentes de software y dependencias,
con información sobre esos componentes y sus relaciones.
\begin{itemize}[leftmargin=1.4em,itemsep=1pt]
  \item Debe ser lo \textbf{más completo posible} e indicar explícitamente dónde no se pueden
        articular las relaciones.
  \item Incluye software de código abierto y con licencia comercial; puede estar ampliamente
        disponible o tener acceso restringido para proteger info confidencial.
  \item Componentes típicos (slide): Libraries, Software Components Name, Author Name, Supplier
        Name, Licences, Version String, Open Source Components, Dependencies, Compliance
        Requirements.
\end{itemize}
\end{defbox}

\begin{anecbox}{La historia del SBOM}
Es una idea \emph{vieja} (un inventario formal de componentes). Se dejó de mantener en cierto
momento, y muchas empresas terminaron \textbf{sin saber qué componentes usaban} (de ahí el
desastre Log4Shell). Hoy \textbf{resurge} por los incidentes recientes; hay plugins en los IDEs
que detectan qué usás, si quedó obsoleto o tiene vulnerabilidades.
\end{anecbox}

%=====================================================================
\section{Protección de la nube}
%=====================================================================

Ya se vio el \textbf{modelo de seguridad compartida} de la nube: la \textbf{configuración es
responsabilidad del cliente}. El uso de \textbf{infraestructura como código (IaC)} permite
definir la configuración de la nube de forma \textbf{reproducible}.

\begin{lstlisting}[style=hcl,caption={Ejemplo de IaC con Terraform (de la slide)},captionpos=b]
resource "aws_instance" "example_server" {
  ami           = "ami-04e914639d0cca79a"
  instance_type = "t2.micro"

  tags = {
    Name = "JacksBlogExample"
  }
}
\end{lstlisting}

\begin{anecbox}{Configuraciones que quedan expuestas}
\prof{Hay redes que quedaron expuestas a la nada, a Internet [\dots] estás exponiendo puertos
que no se están usando o demasiados puertos.} Como IaC ya tiene varios años y está madurando,
aparecen herramientas que detectan \textbf{patrones inseguros} y los marcan. Estas herramientas
están pensadas para empresas con \emph{miles y miles de servidores}; no las usa una Pyme.
\end{anecbox}

\begin{defbox}{Familias de herramientas de protección cloud}
\begin{description}[leftmargin=0pt,labelsep=0.4em,style=nextline,itemsep=3pt]
  \item[CSPM — Cloud Security Posture Management] Busca \textbf{problemas de configuración y
        permisos}. Analiza IaC \emph{antes} de que sea desplegado y revisa la configuración
        actual para ver que no haya habido cambios. Enfoque más \textbf{estático}. (Analiza
        Infrastructure-as-Code y la valida.)
  \item[CWPP — Cloud Workflow/Workload Protection Platform] (en las slides aparece como
        ``CWSS''). Se focaliza en la \textbf{protección de las aplicaciones corriendo} en la
        nube: servidores virtuales, contenedores y \emph{serverless} (ej.\ Lambdas). Analiza
        comportamiento, microsegmentación a nivel de contenedores, threat intelligence.
        \textbf{Complementa a CSPM} (que es más estático).
  \item[CIEM — Cloud Infrastructure Entitlement Management] Foco en los \textbf{permisos de los
        sujetos sobre los recursos} de la nube. Busca \textbf{permisos excesivos}: compara los
        registros de uso de los accesos contra los permisos otorgados para identificar y remover
        los no utilizados. Muestra consistencia entre cuentas cloud y permite una visualización
        efectiva de permisos (útil con múltiples cuentas).
  \item[CNAPP — Cloud Native Application Protection Platform] La ``sombrilla'' que integra todo
        lo anterior. Vista simplificada (Gartner): combina \emph{seguridad proactiva}
        (development artifact scanning) con \emph{reactiva} (runtime visibility/protection),
        atravesada por DevSecOps y \emph{cloud configuration and compliance}.
\end{description}
\end{defbox}

%=====================================================================
\section{Evaluaciones de seguridad}
%=====================================================================

\begin{anecbox}{¿Por qué evaluar la seguridad?}
\prof{Construimos y tratamos de hacer nuestro sistema seguro, tomamos todas las precauciones y
las mejores prácticas\dots y en algún momento ese sistema va a fallar. No pregunten cuándo, no
pregunten cómo, pero va a fallar, así hayamos tomado todos los recaudos. ¿Por qué? Porque la
seguridad no está asegurada.} Van surgiendo nuevos vectores de ataque que uno no contempla; por
eso hay que estar \textbf{evaluando todo el tiempo}: es un proceso constante y sin pausa.
\end{anecbox}

\begin{slidebox}
\textit{``Security doesn't have to last forever; just longer than everything else that might
notice it's gone.''} \hfill (cita de la slide)
\end{slidebox}

\subsection{Tipos de evaluaciones de seguridad}

\begin{defbox}{Las 6 evaluaciones de seguridad}
\begin{description}[leftmargin=0pt,labelsep=0.4em,style=nextline,itemsep=2pt]
  \item[Evaluación de Vulnerabilidades] Uso de \textbf{herramientas automatizadas} para escanear
        sistemas y redes en busca de \emph{vulnerabilidades conocidas}. Da un resumen de
        debilidades explotables.
  \item[Pruebas de Penetración (PenTest)] Más \textbf{exhaustivas}; realizadas por expertos que
        intentan \emph{explotar vulnerabilidades de manera controlada}. Ayudan a entender cómo
        un atacante podría acceder.
  \item[Evaluación de Riesgos] Identificar, evaluar y priorizar riesgos para los activos de
        información. Visión integral (la seguridad como una amenaza más del negocio).
  \item[Auditoría de Seguridad] Revisión \emph{sistemática} de políticas, procedimientos y
        controles. Interna o externa; busca asegurar el cumplimiento de normativas y estándares.
  \item[Evaluación de Cumplimiento] Verifica que la organización cumpla regulaciones y
        estándares aplicables: \textbf{GDPR, HIPAA, PCI-DSS}. Obligatoria para ciertas
        regulaciones.
  \item[Evaluación de Seguridad Física] A menudo se pasa por alto pero es crucial: control de
        acceso a instalaciones, cámaras, protección de archivos físicos y equipos.
  \item[Evaluación de Seguridad de Aplicaciones] Vulnerabilidades en software y apps web:
        revisiones de código, pruebas de seguridad y análisis de arquitectura.
\end{description}
\end{defbox}

\subsection{Penetration Test (PenTest)}

\begin{defbox}{Objetivo del PenTest}
Agarrar el \textbf{sistema funcionando} y probarlo integralmente para \textbf{identificar
vulnerabilidades} que permitan que el sistema pase de un \textbf{estado seguro a uno inseguro}.
\prof{Muchas veces el diseño es bueno y seguro, pero si le pasás ciertos inputs o ante cierto
comportamiento, deja de ser seguro.}
\end{defbox}

\begin{itemize}[leftmargin=1.4em,itemsep=1pt]
  \item Se realizan con un \textbf{alcance definido}, en \textbf{tiempo y sistemas evaluados}
        (podrían extenderse indefinidamente).
  \item Pueden hacerse \textbf{metódicamente}, como control recurrente para verificar que los
        controles definidos siguen en su lugar.
  \item Pueden ejecutarlos un \textbf{grupo de expertos} (internos o un proveedor contratado),
        para identificar problemas no detectados antes.
\end{itemize}

\subsubsection{Tipos de PenTest según visibilidad del evaluador}

\begin{center}
\begin{tabular}{@{}>{\bfseries}p{0.18\textwidth}p{0.76\textwidth}@{}}
\toprule
Tipo & Descripción \\
\midrule
Black Box & El evaluador \emph{no tiene conocimiento previo} del sistema (a lo sumo el dominio o
            IP). Simula un atacante externo sin info interna. También llamado \textbf{Blind Test}. \\
White Box & El evaluador tiene \emph{acceso completo} a la info interna: código fuente, diagramas
            de red, credenciales. Evaluación exhaustiva de vulnerabilidades internas. \\
Gray Box & Enfoque \emph{intermedio}: algún conocimiento, sin acceso completo. Simula un usuario
           con acceso limitado (empleado con privilegios restringidos). \\
Double Blind & Es un ataque Black Box donde \emph{el atacado tampoco está al tanto}. Evalúa
               adicionalmente las capacidades de \textbf{monitoreo y respuesta}. Es lo más
               parecido a la realidad. \\
\bottomrule
\end{tabular}
\end{center}

\begin{anecbox}{¿Por qué Black Box si White Box parece más eficiente?}
Pregunta de un alumno. Respuesta del profe: \prof{El Black Box te da descubrir vectores que vos,
estando sesgado con toda la información, no verías. Estar afuera con la presión de ``tengo que
lograr entrar'' te da un montón de creatividad; con toda la info te metés en la piel del
desarrollador y te terminás encerrando en esa caja sin poder salir.} Además, la razón más clara
para Black Box es que \textbf{no querés compartir} info sensible (red, credenciales, código
fuente) al proveedor, muchas veces por motivos \emph{legales}. \prof{Generalmente se suelen hacer
de los dos tipos.}
\end{anecbox}

\begin{anecbox}{Double Blind y Chaos Monkey}
El \emph{Double Blind} se parece a la metodología \textbf{Chaos Monkey}: ir rompiendo a propósito
aplicaciones dentro de la infraestructura (interrupciones de servicio) \emph{sin avisar} y ver
cómo reacciona el equipo. Si está bien focalizado y controlado, genera aprendizaje sin bajar la
operación.
\end{anecbox}

\subsubsection{Tipos de PenTest según el sistema evaluado}
\begin{itemize}[leftmargin=1.4em,itemsep=1pt]
  \item \textbf{Test de Red Externa}: sistemas y servicios expuestos a Internet (atacantes
        externos).
  \item \textbf{Test de Red Interna}: red interna; simula un ataque desde dentro (empleado
        descontento o malware en una workstation).
  \item \textbf{Test de redes inalámbricas}: NFC, Bluetooth, Wireless, entre otras.
  \item \textbf{Test de Aplicaciones Web}: inyecciones SQL, XSS, etc.
  \item \textbf{Test de Ingeniería Social}: susceptibilidad de los empleados a ser engañados
        (phishing, smishing, campañas con supuestos regalos).
\end{itemize}

\subsubsection{Etapas de un PenTest (6 pasos)}

\begin{center}
\begin{tikzpicture}[font=\scriptsize\bfseries,>=Stealth]
  \def\R{2.7}
  \foreach \ang/\lbl/\i in {90/{1. Preparación}/1, 30/{2. Construir\\plan ataque}/2,
                            -30/{3. Seleccionar\\equipo}/3, -90/{4. Tipo de\\datos a robar}/4,
                            -150/{5. Ejecutar\\el test}/5, 150/{6. Integrar\\resultados}/6}{
    \node[circle,fill=itbaCyan!85,text=black,minimum size=1.35cm,align=center,font=\tiny\bfseries] (n\i) at (\ang:\R) {\lbl};
  }
  \node[align=center,font=\bfseries,itbaCyanDark] at (0,0) {The 6 steps\\of penetration\\testing};
  \foreach \a/\b in {1/2,2/3,3/4,4/5,5/6}{\draw[->,line width=1.2pt,itbaCyanDark] (n\a) to[bend left=12] (n\b);}
  \draw[->,line width=1.2pt,itbaCyanDark] (n6) to[bend left=12] (n1);
\end{tikzpicture}
\end{center}

\begin{enumerate}[leftmargin=1.6em,itemsep=2pt]
  \item \textbf{Preparación.} Identificar el objetivo a evaluar, el alcance y \textbf{acordar el
        nivel de riesgo} que la organización desea correr (ej.: un ataque de fuerza bruta podría
        dejar cuentas bloqueadas y afectar sistemas). Entender legalidades y propósito.
  \item \textbf{Armado del plan de ataque.} Definir herramientas a utilizar y la ubicación
        lógica y física desde donde se lanzará.
  \item \textbf{Identificar el equipo de trabajo.} Expertos en distintas tecnologías y
        generalistas con conocimiento amplio.
  \item \textbf{Definir el objetivo a alcanzar.} Estos ataques son muy enfocados; muchas veces
        \emph{no aportan} info sobre la efectividad del sistema como un todo si no logran el
        objetivo. \prof{Es mejor no tener éxito, porque eso implica que no hay que hacer cambios.}
  \item \textbf{Ejecutar el ataque.} Bajo el tiempo y alcance definidos. Termina al lograr el
        objetivo o al vencerse el tiempo.
  \item \textbf{Integrar los hallazgos} y reiniciar el ciclo, \emph{pivoteando} y escalando con
        lo encontrado. Ejemplo del profe: \emph{(1)} ingeniería social $\to$ \emph{(2)} ataque
        de phishing $\to$ \emph{(3)} con las credenciales se accede al repositorio de archivos y
        se extraen documentos $\to$ con esos documentos se prepara el siguiente ataque.
\end{enumerate}

\subsection{Red / Blue / Purple Team}

Son divisiones dentro del equipo de seguridad de la organización.

\begin{center}
\begin{tikzpicture}[font=\footnotesize]
  \begin{scope}
    \clip (-0.7,0) circle (2.1);
    \fill[attackRed!75] (-2.0,0) circle (2.1);
    \fill[purpleZTA!80] (0.7,0) circle (2.1);
  \end{scope}
  \fill[attackRed!75] (-2.0,0) circle (2.1);
  \fill[blue!65!black] (2.0,0) circle (2.1);
  \begin{scope}
    \clip (-2.0,0) circle (2.1);
    \fill[purpleZTA!85] (2.0,0) circle (2.1);
  \end{scope}
  \node[white,font=\bfseries,align=center] at (-2.5,0.4) {RED TEAM\\\scriptsize OFFENSE};
  \node[white,font=\bfseries,align=center] at (2.5,0.4) {BLUE TEAM\\\scriptsize DEFENSE};
  \node[white,font=\bfseries,align=center] at (0,0.4) {PURPLE\\TEAM};
  \node[white,font=\scriptsize,align=center] at (-2.5,-0.6) {Vuln.\ Assessments\\Pen Tests\\Social Eng.};
  \node[white,font=\scriptsize,align=center] at (2.5,-0.6) {Implem.\ Controls\\Monitoring\\Incident Resp.};
  \node[white,font=\scriptsize,align=center] at (0,-0.5) {Common\\goal};
\end{tikzpicture}
\end{center}

\begin{defbox}{Red / Blue / Purple Team}
\begin{description}[leftmargin=0pt,labelsep=0.4em,style=nextline,itemsep=2pt]
  \item[Red Team (ofensivo)] Piensa y actúa como un \textbf{atacante malintencionado}, usando
        técnicas de hacking para comprometer sistemas, redes y aplicaciones \emph{de forma
        ética y controlada}. Ejecuta pen-tests, ingeniería social y assessments de
        vulnerabilidades, para probar la efectividad de las medidas y mejorar la detección y
        respuesta.
  \item[Blue Team (defensivo)] Implementa y mantiene medidas de seguridad (firewalls, IDS,
        gestión de parches y configuraciones), \textbf{monitorea} en busca de actividad
        sospechosa y \textbf{responde} a intentos de intrusión. Objetivo: proteger los activos y
        mitigar riesgos de ciberataques.
  \item[Purple Team (colaborativo)] Enfoque que \textbf{combina} los esfuerzos del Red y el Blue
        Team para mejorar la postura de seguridad. En lugar de trabajar aislados, fomenta la
        comunicación y colaboración entre ofensiva y defensiva. \prof{Muchas veces surge de la
        pelea entre ambos teams.}
  \end{description}
\end{defbox}

\subsection{Breach and Attack Simulation (BAS)}

\begin{defbox}{Breach and Attack Simulation (BAS)}
Categoría de herramientas y técnicas para \textbf{evaluar la efectividad de las defensas} de una
organización contra posibles ataques. Simulan las \textbf{tácticas, técnicas y procedimientos
(TTP)} que un atacante real podría emplear.
\begin{itemize}[leftmargin=1.4em,itemsep=1pt]
  \item Objetivo: identificar vulnerabilidades y brechas \emph{antes} de que las exploten
        atacantes reales, y medir la \textbf{capacidad de la organización para identificar y
        reaccionar} ante ataques.
  \item Se evalúan las \emph{mismas herramientas} que protegen la organización y los sistemas de
        monitoreo y alertas; identifican cambios de configuración o procesos que dejen de actuar
        como se espera.
  \item Se ejecutan de \textbf{forma recurrente}, probando en distintos puntos la capacidad de
        detección, mitigación y respuesta.
\end{itemize}
\end{defbox}

\begin{anecbox}{``War games'': baratos pero valiosos}
El profe lo asocia a los \emph{juegos de guerra}: se simula un escenario desfavorable puntual,
p.\ ej.\ \prof{¿qué pasa si un atacante obtiene acceso a una ruta/credencial de admin del cloud?
Tiene la ruta, puede borrar infraestructura, crear nueva, hacer lo que quiera.} Se analiza qué
tácticas y técnicas usar para defenderse y qué equipos se involucran. \prof{Son baratos entre
comillas: solamente juntar a la gente involucrada y proponer el escenario.} A veces se prueba
incluso en un ambiente controlado.
\end{anecbox}

%=====================================================================
\section*{Resumen visual: el panorama de la clase}
\addcontentsline{toc}{section}{Resumen visual}
%=====================================================================

\begin{center}
\begin{tikzpicture}[font=\scriptsize,>=Stealth,
   raiz/.style={rectangle,rounded corners,fill=itbaCyanDark,text=white,font=\bfseries\small,minimum height=0.8cm,align=center},
   rama/.style={rectangle,rounded corners,draw=itbaCyanDark,fill=itbaBlue!30,align=center,minimum height=0.6cm,text width=2.7cm}]
  \node[raiz] (r) {Seguridad en la empresa};
  \node[rama,above left=0.7cm and 1.8cm of r] (gov) {Gobernanza};
  \node[rama,left=2.2cm of r] (risk) {Gestión de riesgos\\(Identify/Assess/\\Mitigate/Monitor)};
  \node[rama,below left=0.7cm and 1.8cm of r] (pol) {Políticas\\(Policy$\to$Baselines)};
  \node[rama,above right=0.7cm and 1.8cm of r] (estr) {Estrategias\\(Defense-in-depth\\/ Zero Trust)};
  \node[rama,right=2.2cm of r] (pam) {PAM + DevSecOps\\(IAM, Shift Left)};
  \node[rama,below right=0.7cm and 1.8cm of r] (eval) {Evaluaciones\\(PenTest, Red/Blue,\\BAS)};
  \foreach \n in {gov,risk,pol,estr,pam,eval}{\draw[->,itbaCyanDark,line width=1pt] (r) -- (\n);}
\end{tikzpicture}
\end{center}

%=====================================================================
\section*{Lectura recomendada}
\addcontentsline{toc}{section}{Lectura recomendada}
%=====================================================================
\begin{center}
\begin{tcolorbox}[enhanced,colback=itbaBlue!20,colframe=itbaCyanDark,boxrule=1pt,
  arc=3pt,width=0.85\textwidth,halign=center]
\large\textbf{NIST SP 800-207 — Zero Trust Architecture}\\[4pt]
Capítulos 2 y 3\\[2pt]
\url{https://csrc.nist.gov/pubs/sp/800/207/final}
\end{tcolorbox}
\end{center}
\noindent\textbf{Nota de alcance.} El profe aclaró que \emph{esta clase no está en Bishop}: el
libro tiene material relacionado, pero esta clase está enfocada en el panorama profesional
contemporáneo y se apoya en el estándar \textbf{NIST 800-207}. El detalle del \emph{Penetration
Testing} (metodología paso a paso) corresponde además a la clase aparte de \textbf{Pentesting
(Clase 10b)}.

%---------------------------------------------------------------------
\vfill
\begin{center}\small\itshape\color{black!60}
Apunte integral de la Clase 10 — Seguridad en la empresa · Criptografía y Seguridad, ITBA.\\
Síntesis de teoría (transcripciones + diapositivas) y práctica.
\end{center}

\end{document}
